\section{Experiments}

\subsection{Data and evaluation}
We evaluated online identity reconnection using the native manual keyframes of
PersonPath22 \cite{shuai2022personpath22}. We used the publisher's FCOS
person detections, trained on CrowdHuman, and aligned each detection frame
with the corresponding annotated source frame. The resulting observation grid
was nominally 5 Hz. Of the 236 public clips, we assigned 16 publisher-hosted
training assets to development and 91 to calibration. We retained 74
publisher-hosted official test assets for evaluation. We kept the 55 clips
with external source mappings separate from this protocol. The detection
input, annotation matching rules, and evaluation sequences were identical
across tracker variants.

To measure detector-gap recovery, we first matched detection boxes and
ground-truth persons one to one at IoU $\geq 0.50$. For each ground-truth
identity, a gap consisted of one to five intervening annotated keyframes with
no matched detection. We counted a recovery if the track identifier at the
first returned detection matched its identifier before the gap. We used the
same ground-truth/detector gap list for all policies and included cases
without an emitted pre-gap identity in the denominator.

We also audited the tracker's reactivation actions separately from gap
recovery. A reactivation was correct if the specified one-to-one matching
assigned its previous emitted track and its new detector box to the same
annotated person. Otherwise
it was counted as wrong or unverifiable, except where an unmatched endpoint
overlapped an annotated ignore region by more than half its own box area.
We report these ignored actions separately. This evaluation is performed
after tracking. The gate sees the detector and track state alone. Overall
metrics were HOTA, AssA, IDF1, and ID switches, computed with TrackEval on
the native keyframe grid \cite{luiten2021hota}. PersonPath22's distractor
and crowd treatment was used for the overall metrics.

\subsection{Comparison and model fitting}
We replayed ByteTrack \cite{zhang2022bytetrack} and a motion-only BoT-SORT
configuration \cite{aharon2022botsort}. Each gate accepted or declined a
native old-ID reactivation before the tracker updated its state. Each
policy started from frame one. We compared the unrestricted tracker, a
source cap alone, a threshold on the native association cost, a threshold
on the nearest competitor margin, an additive linear score, our interaction
score, and the interaction score with a within-source ranking term.
All capped policies allowed at most 12 admissions per source before
evaluation. The cap applied to all admissions. The score was fit on the
development sources for 200 updates with Adam (learning rate 0.04 and
coefficient penalty 0.002) on a GeForce RTX 3060 Ti. Candidate thresholds
were fixed from development score quantiles before calibration. We chose
the strongest recovery operating point for each score within a common
empirical calibration risk budget, with at least 20\% of calibration
sources active. We then applied the selected policies to the official
test assets without tuning.

\subsection{Finite-panel risk analysis}
We fixed a second analysis over the 165 non-development publisher-hosted
assets and drew 91 sources uniformly without replacement. We computed
simultaneous one-sided bounds for the frozen ByteTrack family by inverting
the fixed betting-mixture tests. The event-weighted and active-video
estimands each received half of a familywise error budget of 0.05.
The target error budget was 20\%. We also assessed whether the same upper
bounds reached 10\%. This randomization analysis concerns the fixed panel.
The official-test comparison uses strategies chosen on the original
calibration split with a separate selector. The randomized calibration
sample included 41 official-test sources, and its complement included the
other 33. The two analyses use overlapping data and answer different
questions.

\subsection{Official split results}
The official test subset contained 9,320 keyframes and 3,155 eligible
detector gaps. At a common empirical calibration ceiling of 47.23\%,
the ByteTrack interaction gate reduced the tested incorrect or
unverifiable action proportion from 46.77\% to 35.39\% relative to
the native association threshold. It increased verified correct
reactivations from 239 to 283, using 438 evaluated admissions instead
of 449. The gap recoveries were 393 and 388. The paired source
resampling interval was $[-16.57,-6.40]$ percentage points for the
failure-rate difference and $[-0.33,0.67]$ for the gap-recovery
difference (10,000 replicates, comparison-wise). The interaction score
also reduced the failure proportion against the cap-only and linear
controls at this ceiling.

At the 20\% calibration ceiling, both ByteTrack scores recovered 383
gaps, with observed failure proportions of 21.24\% for the interaction
score and 21.15\% for the linear score. On motion-only BoT-SORT, the
interaction score recovered 334 gaps versus 327 for the linear score
and yielded 16.19\% versus 18.31\% checked failures. The paired
resampling intervals were $[0.05,0.48]$ percentage points for gap
recovery and $[-4.94,0.75]$ for the failure proportion. At the
relaxed BoT-SORT ceiling (42.15\%), the linear score had lower
observed failure than the interaction score (25.19\% versus 37.28\%).
The added pairwise ranking loss did not produce a consistent benefit.

\begin{table*}[t]
\centering
\small
\caption{PersonPath22 native-keyframe official-test subset. Gap recovery
uses the same 3,155 gaps for each policy. The evaluated admits and
fails columns exclude ignored actions. HOTA, AssA and IDF1 are on a
0--100 scale, and IDSW is a count. The risk column gives the observed
failure proportion. The cap-only gate applies the same source cap.
Within each tracker, rows without a strict marker use the relaxed
calibration ceiling for score selection.}
\label{tab:reconnection-main}
\begin{tabular}{llrrrrrrrr}
\hline
Tracker & Reactivation policy & Gaps kept & Admits & Fails & Risk (\%) & HOTA & AssA & IDF1 & IDSW \\
\hline
OC-SORT & External reference & 174 & -- & -- & -- & 41.83 & 41.22 & 52.15 & 3790 \\
\hline
ByteTrack & Unrestricted & 529 & 2589 & 1552 & 59.95 & 46.64 & 42.33 & 56.73 & 6662 \\
 & Cap only & 399 & 537 & 310 & 57.73 & 44.69 & 39.32 & 52.94 & 7129 \\
 & Native cost & 388 & 449 & 210 & 46.77 & 44.76 & 39.41 & 53.16 & 7104 \\
 & Nearest margin & 391 & 456 & 202 & 44.30 & 44.84 & 39.57 & 53.14 & 7132 \\
 & Linear & 395 & 514 & 247 & 48.05 & 44.90 & 39.63 & 53.38 & 7131 \\
 & Interaction & 393 & 438 & 155 & 35.39 & 44.99 & 39.81 & 53.38 & 7022 \\
 & Interaction + ranking & 393 & 509 & 240 & 47.15 & 44.91 & 39.69 & 53.34 & 7131 \\
 & Linear (strict) & 383 & 331 & 70 & 21.15 & 44.87 & 39.63 & 53.25 & 7120 \\
 & Interaction (strict) & 383 & 339 & 72 & 21.24 & 44.87 & 39.63 & 53.23 & 7116 \\
\hline
BoT-SORT & Unrestricted & 424 & 1858 & 992 & 53.39 & 46.84 & 43.78 & 56.57 & 5909 \\
(motion only) & Cap only & 354 & 495 & 241 & 48.69 & 45.04 & 40.71 & 53.03 & 6616 \\
 & Native cost & 343 & 392 & 148 & 37.76 & 44.99 & 40.67 & 52.98 & 6656 \\
 & Nearest margin & 347 & 451 & 185 & 41.02 & 45.10 & 40.83 & 53.07 & 6654 \\
 & Linear & 339 & 397 & 100 & 25.19 & 45.20 & 40.99 & 53.29 & 6603 \\
 & Interaction & 346 & 448 & 167 & 37.28 & 45.14 & 40.88 & 53.20 & 6566 \\
 & Interaction + ranking & 344 & 450 & 170 & 37.78 & 45.14 & 40.87 & 53.16 & 6611 \\
 & Linear (strict) & 327 & 284 & 52 & 18.31 & 45.05 & 40.83 & 53.05 & 6628 \\
 & Interaction (strict) & 334 & 278 & 45 & 16.19 & 45.06 & 40.85 & 53.02 & 6626 \\
\hline
\end{tabular}
\end{table*}

The unrestricted trackers retained more natural gaps and attained
higher overall HOTA as shown in Table~\ref{tab:reconnection-main}.
Relative to unrestricted ByteTrack, the relaxed interaction gate
reduced evaluated failures by 1,397 and retained 136 fewer natural gaps.
HOTA declined by 1.65 points. These changes reflect the cost of selective
admission. We include OC-SORT \cite{cao2023ocsort} as an overall
tracking reference. Its output supported evaluating natural-gap
continuity, while its internal reactivation mechanism used a different
protocol.

\subsection{Randomized finite panel}
Two of the 23 frozen ByteTrack policies met the 20\% active-video
bound. The interaction policy had an upper bound of 19.95\% and the
linear policy had an upper bound of 19.72\%. The interaction policy
was selected because it recovered 752 calibration gaps compared
with 732 for the linear policy. None of the evaluated policies met
the event-weighted 20\% bound. Its lowest upper bound was 32.28\%.
The two probability budgets each used 0.025 familywise error.

On the 74-source random complement, the selected interaction policy
accepted 192 evaluated reactivations, of which 16 failed, and was
active on 43 videos. The observed active-video mean was 9.34\%.
It recovered 359 of 2,830 gaps. The eligible linear policy accepted
74 actions, of which seven failed, with 353 recovered gaps. It was
active on 26 videos. These observations describe the held-out complement.

\subsection{Transfer across datasets}
We transferred the PersonPath22 ByteTrack thresholds to 40 previously
unseen clips under the DanceTrack training partition, with a frozen
COCO YOLOX-X detector. The relaxed interaction gate recovered 1,506 of
8,472 gaps and failed in 199 of 462 accepted actions (43.07\%). The
native-cost gate recovered 1,500 gaps and failed in 212 of 463 actions
(45.79\%). The 20\%-calibration interaction gate had 30.03\% evaluated
failures in DanceTrack. In a five-video exploratory BDD100K transfer,
the ByteTrack interaction gate at the stricter operating point had
three failures among eight accepted actions, compared with 16 among
25 for unrestricted ByteTrack. Each policy retained five of the 121
eligible detector gaps. These data characterize transfer behavior
under new detection inputs and acquisition conditions.

\subsection{Label sensitivity and execution cost}
We reassigned numerically equivalent duplicate detector boxes
within 0.001 pixels per coordinate when their matched ground-truth
identity was unambiguous. This post hoc check changed none of the
accepted or failed action counts in the relaxed ByteTrack
native-cost and interaction comparisons or the strict BoT-SORT
linear and interaction comparisons. The relaxed ByteTrack
interaction gate retained annotated identity 66 across keyframes
47--49 of \texttt{uid\_vid\_00013}, where the native threshold
missed the return. The converse occurred for identity 20 across
keyframes 66--68 of \texttt{uid\_vid\_00008}.

Median association and decision-logging time was 1.60 and
1.54 ms per keyframe for unrestricted and relaxed interaction
ByteTrack, and 1.48 and 1.49 ms for unrestricted and strict
interaction BoT-SORT. These timings cover seven CPU replays, each
with 293 development keyframes, excluding detector inference and
evaluation.
